\documentclass[a4paper,12pt]{article}
\usepackage[most]{tcolorbox}
\usepackage{xcolor}

\title{Algos condensed and explained}
\author{salimek}
\date{\today}

\begin{document}

\maketitle



\section{Pattern}

\begin{tcolorbox}[
    colback=gray!30,
    colframe=black,
    boxrule=1pt,
    arc=0mm
]

\begin{tcolorbox}[
    colback=gray!10,
    colframe=black,
    arc=0mm
]
\centering
{\Huge\bfseries TITLE}
\end{tcolorbox}

\begin{tcolorbox}[
    colback=gray!10,
    colframe=black,
    arc=0mm
]
\small
Explanation of the algorithm.
\end{tcolorbox}

\begin{tcolorbox}[
    colback=black!5,
    colframe=black,
    arc=0mm,
    title=Pseudo-code
]
\small
\ttfamily
Algorithm(...)\\
\quad instruction 1\\
\quad instruction 2\\
\quad return result
\end{tcolorbox}

\begin{tcolorbox}[
    colback=cyan!20,
    colframe=black,
    arc=0mm
]
\small
\textbf{Time Complexity:} $O(n)$
\end{tcolorbox}

\begin{tcolorbox}[
    colback=violet!20,
    colframe=black,
    arc=0mm
]
\small
\textbf{Space Complexity:} $O(1)$
\end{tcolorbox}

\begin{tcolorbox}[
    colback=pink!20,
    colframe=black,
    arc=0mm
]
\small
\textbf{Recurrence Formula:}

\[
T(n)=2T\left(\frac{n}{2}\right)+n
\]
\end{tcolorbox}

\end{tcolorbox}























\section{Algos (everything may not be finished and explained)}


\begin{tcolorbox}[
    colback=gray!30,
    colframe=black,
    boxrule=1pt,
    arc=0mm
]

\begin{tcolorbox}[
    colback=gray!10,
    colframe=black,
    arc=0mm
]
\centering
{\Huge\bfseries INSERTION SORT}
\end{tcolorbox}

\begin{tcolorbox}[
    colback=gray!10,
    colframe=black,
    arc=0mm
]
\small
Insertion Sort is a simple sorting algorithm that builds the final sorted array one element at a time. 
At each iteration, the current element is inserted into its correct position inside the already sorted part of the array.
\end{tcolorbox}

\begin{tcolorbox}[
    colback=black!5,
    colframe=black,
    arc=0mm,
    title=Pseudo-code
]
\small
\begin{lstlisting}
Insertion-Sort(A, n)

    for j = 2 to n
        key = A[j]
        i = j - 1

        while i > 0 and A[i] > key
            A[i + 1] = A[i]
            i = i - 1

        A[i + 1] = key
\end{lstlisting}
\end{tcolorbox}

\begin{tcolorbox}[
    colback=cyan!20,
    colframe=black,
    arc=0mm
]
\small
\textbf{Time Complexity:} $\Theta(n^2)$ worst/average, $\Theta(n)$ best case
\end{tcolorbox}

\begin{tcolorbox}[
    colback=violet!20,
    colframe=black,
    arc=0mm
]
\small
\textbf{Space Complexity:} $O(1)$
\end{tcolorbox}

\begin{tcolorbox}[
    colback=pink!20,
    colframe=black,
    arc=0mm
]
\small
\textbf{Recurrence Formula:} N/A
\end{tcolorbox}

\end{tcolorbox}







\begin{tcolorbox}[
    colback=gray!30,
    colframe=black,
    boxrule=1pt,
    arc=0mm
]

\begin{tcolorbox}[
    colback=gray!10,
    colframe=black,
    arc=0mm
]
\centering
{\Huge\bfseries Merge Sort}
\end{tcolorbox}

\begin{tcolorbox}[
    colback=gray!10,
    colframe=black,
    arc=0mm
]
\small
Divide the array into two halves, recursively sort each half, and merge the sorted halves.
The objective is to sort an array efficiently using the divide-and-conquer strategy.
\end{tcolorbox}

\begin{tcolorbox}[
    colback=black!5,
    colframe=black,
    arc=0mm,
    title=Pseudo-code
]
\tiny
\begin{lstlisting}
MERGE-SORT(A, p, r)

    if p < r
        q = rounddown((p+r)/2)

        MERGE-SORT(A, p, q)
        MERGE-SORT(A, q+1, r)

        MERGE(A, p, q, r)


MERGE(A,p,q,r)

    n1 = q-p+1
    n2 = r-q

    create arrays L[1...n1+1], R[1...n2+1]

    for i = 1 to n1
        L[i] = A[p+i-1]

    for j = 1 to n2
        R[j] = A[q+j]

    L[n1+1] = 8
    R[n2+1] = 8

    i = 1
    j = 1

    for k = p to r
        if L[i] <= R[j]
            A[k] = L[i]
            i = i + 1

        else
            A[k] = R[j]
            j = j + 1
\end{lstlisting}
\end{tcolorbox}

\begin{tcolorbox}[
    colback=cyan!20,
    colframe=black,
    arc=0mm
]
\small
\textbf{Time Complexity:} $\Theta(n \log n)$ in all cases.
\end{tcolorbox}

\begin{tcolorbox}[
    colback=violet!20,
    colframe=black,
    arc=0mm
]
\small
\textbf{Space Complexity:} $O(n)$ auxiliary space.
\end{tcolorbox}

\begin{tcolorbox}[
    colback=pink!20,
    colframe=black,
    arc=0mm
]
\small
\textbf{Recurrence Formula:}
$T(n)=2T(n/2)+\Theta(n)$
$\Rightarrow T(n)=\Theta(n\log n)$ by Master Theorem (case 2).
\end{tcolorbox}

\end{tcolorbox}


















































\begin{tcolorbox}[
    colback=gray!30,
    colframe=black,
    boxrule=1pt,
    arc=0mm
]

\begin{tcolorbox}[
    colback=gray!10,
    colframe=black,
    arc=0mm
]
\centering
{\Huge\bfseries Maximum Subarray Optimized}
\end{tcolorbox}


\begin{tcolorbox}[
    colback=gray!10,
    colframe=black,
    arc=0mm
]
\small
Find the continuous subarray with the maximum possible sum.
The algorithm uses divide and conquer by splitting the array,
solving both halves recursively, and checking the subarray crossing the middle.
\end{tcolorbox}


\begin{tcolorbox}[
    colback=black!5,
    colframe=black,
    arc=0mm,
    title=Pseudo-code
]
\tiny
\begin{lstlisting}
FIND-MAXIMUM-SUBARRAY(A, low, high)

    if high == low
        return (low, high, A[low])

    mid = (low + high) / 2

    left  = FIND-MAXIMUM-SUBARRAY(A, low, mid)
    right = FIND-MAXIMUM-SUBARRAY(A, mid+1, high)

    cross = FIND-MAX-CROSSING(A, low, mid, high)

    if left-sum >= right-sum and left-sum >= cross-sum
        return (left-low, left-high, left-sum)

    elseif right-sum >= left-sum and right-sum >= cross-sum
        return (right-low, right-high, right-sum)

    else
        return (cross-low, cross-high, cross-sum)



FIND-MAX-CROSSING(A, low, mid, high)

    left-sum = -8
    sum = 0

    for i = mid downto low
        sum = sum + A[i]

        if sum > left-sum
            left-sum = sum
            max-left = i


    right-sum = -8
    sum = 0

    for j = mid+1 to high
        sum = sum + A[j]

        if sum > right-sum
            right-sum = sum
            max-right = j


    return (max-left, max-right, left-sum + right-sum)
\end{lstlisting}
\end{tcolorbox}


\begin{tcolorbox}[
    colback=cyan!20,
    colframe=black,
    arc=0mm
]
\small
\textbf{Time Complexity:}
$T(n)=2T(n/2)+\Theta(n)$
$\Rightarrow \Theta(n\log n)$ by Master Theorem.
\end{tcolorbox}


\begin{tcolorbox}[
    colback=violet!20,
    colframe=black,
    arc=0mm
]
\small
\textbf{Space Complexity:}
$O(\log n)$ auxiliary space due to recursion stack.
\end{tcolorbox}


\begin{tcolorbox}[
    colback=pink!20,
    colframe=black,
    arc=0mm
]
\small
\textbf{Recurrence Formula:}
$T(n)=2T(n/2)+\Theta(n)$
\end{tcolorbox}

\end{tcolorbox}










































\begin{tcolorbox}[
    colback=gray!30,
    colframe=black,
    boxrule=1pt,
    arc=0mm
]

\begin{tcolorbox}[
    colback=gray!10,
    colframe=black,
    arc=0mm
]
\centering
{\Huge\bfseries Maximum Subarray Brute Force}
\end{tcolorbox}


\begin{tcolorbox}[
    colback=gray!10,
    colframe=black,
    arc=0mm
]
\small
Checks every possible continuous subarray and keeps the one with the largest sum.
The algorithm uses an exhaustive search approach.
\end{tcolorbox}


\begin{tcolorbox}[
    colback=black!5,
    colframe=black,
    arc=0mm,
    title=Pseudo-code
]
\tiny
\begin{lstlisting}
MAXIMUM-SUBARRAY-SLOW(A[1..n])

    B.val = 8
    B.i = 1
    B.j = n

    for i = 1 to n

        tmp = 0

        for j = i to n

            tmp = tmp + A[j]

            if tmp > B.val
                B.val = tmp
                B.i = i
                B.j = j

    return (B.i, B.j, B.val)
\end{lstlisting}
\end{tcolorbox}


\begin{tcolorbox}[
    colback=cyan!20,
    colframe=black,
    arc=0mm
]
\small
\textbf{Time Complexity:}
$\Theta(n^2)$
\end{tcolorbox}


\begin{tcolorbox}[
    colback=violet!20,
    colframe=black,
    arc=0mm
]
\small
\textbf{Space Complexity:}
$\Theta(n)$ input space + $\Theta(1)$ auxiliary space.
\end{tcolorbox}


\begin{tcolorbox}[
    colback=pink!20,
    colframe=black,
    arc=0mm
]
\small
\textbf{Recurrence Formula:}
No recurrence formula (iterative brute force algorithm).
\end{tcolorbox}

\end{tcolorbox}




























\begin{tcolorbox}[
    colback=gray!30,
    colframe=black,
    boxrule=1pt,
    arc=0mm
]

\begin{tcolorbox}[
    colback=gray!10,
    colframe=black,
    arc=0mm
]
\centering
{\Huge\bfseries Matrix Multiplication Naive}
\end{tcolorbox}


\begin{tcolorbox}[
    colback=gray!10,
    colframe=black,
    arc=0mm
]
\small
Multiply two square matrices using the standard triple nested loop approach.
Each element of the result matrix is computed by the dot product of a row and a column.
\end{tcolorbox}


\begin{tcolorbox}[
    colback=black!5,
    colframe=black,
    arc=0mm,
    title=Pseudo-code
]
\tiny
\begin{lstlisting}
SQUARE-MAT-MULT(A, B, n)

    create matrix C of size n x n

    for i = 1 to n

        for j = 1 to n

            C[i][j] = 0

            for k = 1 to n

                C[i][j] = C[i][j] + A[i][k] * B[k][j]

    return C
\end{lstlisting}
\end{tcolorbox}


\begin{tcolorbox}[
    colback=cyan!20,
    colframe=black,
    arc=0mm
]
\small
\textbf{Time Complexity:}
$O(n^3)$
\end{tcolorbox}


\begin{tcolorbox}[
    colback=violet!20,
    colframe=black,
    arc=0mm
]
\small
\textbf{Space Complexity:}
$O(n^2)$ for the result matrix.
\end{tcolorbox}


\begin{tcolorbox}[
    colback=pink!20,
    colframe=black,
    arc=0mm
]
\small
\textbf{Recurrence Formula:}
$T(n)=8T(n/2)+\Theta(n^2)$
$\Rightarrow T(n)=\Theta(n^3)$ by Master Theorem (case 1).
\end{tcolorbox}

\end{tcolorbox}









































\begin{tcolorbox}[
    colback=gray!30,
    colframe=black,
    boxrule=1pt,
    arc=0mm
]

\begin{tcolorbox}[
    colback=gray!10,
    colframe=black,
    arc=0mm
]
\centering
{\Huge\bfseries Recursive Matrix Multiplication}
\end{tcolorbox}


\begin{tcolorbox}[
    colback=gray!10,
    colframe=black,
    arc=0mm
]
\small
Multiply two square matrices using a divide-and-conquer approach.
The matrices are divided into smaller submatrices and recursively multiplied.
\end{tcolorbox}


\begin{tcolorbox}[
    colback=black!5,
    colframe=black,
    arc=0mm,
    title=Pseudo-code
]
\tiny
\begin{lstlisting}
REC-MAT-MULT(A, B, n)

    create matrix C of size n x n

    if n == 1

        C11 = A11 * B11


    else

        partition A, B, C into n/2 x n/2 blocks


        C11 = REC-MAT-MULT(A11,B11,n/2)
              + REC-MAT-MULT(A12,B21,n/2)

        C12 = REC-MAT-MULT(A11,B12,n/2)
              + REC-MAT-MULT(A12,B22,n/2)

        C21 = REC-MAT-MULT(A21,B11,n/2)
              + REC-MAT-MULT(A22,B21,n/2)

        C22 = REC-MAT-MULT(A21,B12,n/2)
              + REC-MAT-MULT(A22,B22,n/2)


    return C
\end{lstlisting}
\end{tcolorbox}


\begin{tcolorbox}[
    colback=cyan!20,
    colframe=black,
    arc=0mm
]
\small
\textbf{Time Complexity:}
$T(n)=8T(n/2)+\Theta(n^2)$
$\Rightarrow O(n^3)$ by Master Theorem.
\end{tcolorbox}


\begin{tcolorbox}[
    colback=violet!20,
    colframe=black,
    arc=0mm
]
\small
\textbf{Space Complexity:}
$O(n^2)$ for storing matrices.
\end{tcolorbox}


\begin{tcolorbox}[
    colback=pink!20,
    colframe=black,
    arc=0mm
]
\small
\textbf{Recurrence Formula:}
$T(n)=8T(n/2)+\Theta(n^2)$
\end{tcolorbox}

\end{tcolorbox}































\begin{tcolorbox}[
    colback=gray!30,
    colframe=black,
    boxrule=1pt,
    arc=0mm
]

\begin{tcolorbox}[
    colback=gray!10,
    colframe=black,
    arc=0mm
]
\centering
{\Huge\bfseries Strassen's Algorithm}
\end{tcolorbox}


\begin{tcolorbox}[
    colback=gray!10,
    colframe=black,
    arc=0mm
]
\small
Multiply two matrices using divide and conquer while reducing the number of recursive
multiplications from 8 to 7. The algorithm uses linear combinations of submatrices
to improve the complexity.
\end{tcolorbox}


\begin{tcolorbox}[
    colback=black!5,
    colframe=black,
    arc=0mm,
    title=Pseudo-code
]
\tiny
\begin{lstlisting}
STRASSEN-MAT-MULT(A, B, n)

    create matrix C of size n x n

    if n == 1

        C11 = A11 * B11


    else

        partition A, B, C into n/2 x n/2 blocks


        M1 = (A11 + A22) * (B11 + B22)

        M2 = (A21 + A22) * B11

        M3 = A11 * (B12 - B22)

        M4 = A22 * (B21 - B11)

        M5 = (A11 + A12) * B22

        M6 = (A21 - A11) * (B11 + B12)

        M7 = (A12 - A22) * (B21 + B22)


        C11 = M1 + M4 - M5 + M7

        C12 = M3 + M5

        C21 = M2 + M4

        C22 = M1 - M2 + M3 + M6


    return C
\end{lstlisting}
\end{tcolorbox}


\begin{tcolorbox}[
    colback=cyan!20,
    colframe=black,
    arc=0mm
]
\small
\textbf{Time Complexity:}
$T(n)=7T(n/2)+\Theta(n^2)$
$\Rightarrow \Theta(n^{\log_2 7}) \approx \Theta(n^{2.807})$.
\end{tcolorbox}


\begin{tcolorbox}[
    colback=violet!20,
    colframe=black,
    arc=0mm
]
\small
\textbf{Space Complexity:}
$O(n^2)$ for storing matrices.
\end{tcolorbox}


\begin{tcolorbox}[
    colback=pink!20,
    colframe=black,
    arc=0mm
]
\small
\textbf{Recurrence Formula:}
$T(n)=7T(n/2)+\Theta(n^2)$
\end{tcolorbox}

\end{tcolorbox}


















\begin{tcolorbox}[
    colback=gray!30,
    colframe=black,
    boxrule=1pt,
    arc=0mm
]

\begin{tcolorbox}[
    colback=gray!10,
    colframe=black,
    arc=0mm
]
\centering
{\Huge\bfseries Karatsuba Multiplication}
\end{tcolorbox}


\begin{tcolorbox}[
    colback=gray!10,
    colframe=black,
    arc=0mm
]
\small
Multiply two large integers faster than the classical multiplication method.
The algorithm reduces the number of recursive multiplications from 4 to 3
using a divide-and-conquer strategy.
\end{tcolorbox}


\begin{tcolorbox}[
    colback=black!5,
    colframe=black,
    arc=0mm,
    title=Pseudo-code
]
\tiny
\begin{lstlisting}
KARATSUBA-MULT(x, y, n)

    if n == 1

        return x * y


    else

        split x into xH, xL

        split y into yH, yL


        p1 = KARATSUBA-MULT(xH, yH, n/2)

        p2 = KARATSUBA-MULT(xL, yL, n/2)

        p3 = KARATSUBA-MULT(
             xH + xL,
             yH + yL,
             n/2)


        return p1*b^n
               + (p3 - p1 - p2)*b^(n/2)
               + p2
\end{lstlisting}
\end{tcolorbox}


\begin{tcolorbox}[
    colback=cyan!20,
    colframe=black,
    arc=0mm
]
\small
\textbf{Time Complexity:}
$T(n)=3T(n/2)+\Theta(n)$
$\Rightarrow \Theta(n^{\log_2 3}) \approx \Theta(n^{1.585})$.
\end{tcolorbox}


\begin{tcolorbox}[
    colback=violet!20,
    colframe=black,
    arc=0mm
]
\small
\textbf{Space Complexity:}
$O(1)$ auxiliary space.
\end{tcolorbox}


\begin{tcolorbox}[
    colback=pink!20,
    colframe=black,
    arc=0mm
]
\small
\textbf{Recurrence Formula:}
$T(n)=3T(n/2)+\Theta(n)$
\end{tcolorbox}

\end{tcolorbox}
























\begin{tcolorbox}[
    colback=gray!30,
    colframe=black,
    boxrule=1pt,
    arc=0mm
]

\begin{tcolorbox}[
    colback=gray!10,
    colframe=black,
    arc=0mm
]
\centering
{\Huge\bfseries MAX-HEAPIFY}
\end{tcolorbox}


\begin{tcolorbox}[
    colback=gray!10,
    colframe=black,
    arc=0mm
]
\small
Restore the max-heap property at a given node by comparing it with its children.
If a child is larger, swap values and continue recursively down the heap.
\end{tcolorbox}


\begin{tcolorbox}[
    colback=black!5,
    colframe=black,
    arc=0mm,
    title=Pseudo-code
]
\tiny
\begin{lstlisting}
MAX-HEAPIFY(A, i, n)

    l = LEFT(i)

    r = RIGHT(i)


    if l <= n and A[l] > A[i]

        largest = l

    else

        largest = i



    if r <= n and A[r] > A[largest]

        largest = r



    if largest != i

        exchange A[i] with A[largest]

        MAX-HEAPIFY(A, largest, n)
\end{lstlisting}
\end{tcolorbox}


\begin{tcolorbox}[
    colback=cyan!20,
    colframe=black,
    arc=0mm
]
\small
\textbf{Time Complexity:}
$\Theta(height(i)) = O(\log n)$
\end{tcolorbox}


\begin{tcolorbox}[
    colback=violet!20,
    colframe=black,
    arc=0mm
]
\small
\textbf{Space Complexity:}
$O(\log n)$ due to recursion stack.
\end{tcolorbox}


\begin{tcolorbox}[
    colback=pink!20,
    colframe=black,
    arc=0mm
]
\small
\textbf{Recurrence Formula:}
$T(n)=T(n/2)+\Theta(1)$
$\Rightarrow T(n)=O(\log n)$
\end{tcolorbox}

\end{tcolorbox}




















\begin{tcolorbox}[
    colback=gray!30,
    colframe=black,
    boxrule=1pt,
    arc=0mm
]

\begin{tcolorbox}[
    colback=gray!10,
    colframe=black,
    arc=0mm
]
\centering
{\Huge\bfseries Build Heap}
\end{tcolorbox}


\begin{tcolorbox}[
    colback=gray!10,
    colframe=black,
    arc=0mm
]
\small
Transform an unordered array into a valid max-heap.
The algorithm applies MAX-HEAPIFY bottom-up on all internal nodes
to restore the heap property efficiently.
\end{tcolorbox}


\begin{tcolorbox}[
    colback=black!5,
    colframe=black,
    arc=0mm,
    title=Pseudo-code
]
\tiny
\begin{lstlisting}
BUILD-MAX-HEAP(A, n)

    for i = n/2 downto 1

        MAX-HEAPIFY(A, i, n)
\end{lstlisting}
\end{tcolorbox}


\begin{tcolorbox}[
    colback=cyan!20,
    colframe=black,
    arc=0mm
]
\small
\textbf{Time Complexity:}
$O(n)$
\end{tcolorbox}


\begin{tcolorbox}[
    colback=violet!20,
    colframe=black,
    arc=0mm
]
\small
\textbf{Space Complexity:}
$O(1)$ auxiliary space.
\end{tcolorbox}


\begin{tcolorbox}[
    colback=pink!20,
    colframe=black,
    arc=0mm
]
\small
\textbf{Recurrence Formula:}
No recurrence formula (bottom-up iterative algorithm).
\end{tcolorbox}

\end{tcolorbox}

















\begin{tcolorbox}[
    colback=gray!30,
    colframe=black,
    boxrule=1pt,
    arc=0mm
]

\begin{tcolorbox}[
    colback=gray!10,
    colframe=black,
    arc=0mm
]
\centering
{\Huge\bfseries Heapsort}
\end{tcolorbox}


\begin{tcolorbox}[
    colback=gray!10,
    colframe=black,
    arc=0mm
]
\small
Sort an array using a max-heap structure.
The algorithm first builds a max-heap, then repeatedly extracts
the maximum element and restores the heap property.
\end{tcolorbox}


\begin{tcolorbox}[
    colback=black!5,
    colframe=black,
    arc=0mm,
    title=Pseudo-code
]
\tiny
\begin{lstlisting}
HEAPSORT(A, n)

    BUILD-MAX-HEAP(A, n)


    for i = n downto 2

        exchange A[1] with A[i]

        MAX-HEAPIFY(A, 1, i-1)
\end{lstlisting}
\end{tcolorbox}


\begin{tcolorbox}[
    colback=cyan!20,
    colframe=black,
    arc=0mm
]
\small
\textbf{Time Complexity:}
$O(n \log n)$
\end{tcolorbox}


\begin{tcolorbox}[
    colback=violet!20,
    colframe=black,
    arc=0mm
]
\small
\textbf{Space Complexity:}
$O(1)$ auxiliary space.
\end{tcolorbox}


\begin{tcolorbox}[
    colback=pink!20,
    colframe=black,
    arc=0mm
]
\small
\textbf{Recurrence Formula:}
No recurrence formula (iterative heap-based algorithm).
\end{tcolorbox}

\end{tcolorbox}
























\begin{tcolorbox}[
    colback=gray!30,
    colframe=black,
    boxrule=1pt,
    arc=0mm
]

\begin{tcolorbox}[
    colback=gray!10,
    colframe=black,
    arc=0mm
]
\centering
{\Huge\bfseries Heap Extract Max}
\end{tcolorbox}


\begin{tcolorbox}[
    colback=gray!10,
    colframe=black,
    arc=0mm
]
\small
Remove and return the maximum element of a max-heap.
The root is replaced by the last element, then MAX-HEAPIFY
is used to restore the heap property.
\end{tcolorbox}


\begin{tcolorbox}[
    colback=black!5,
    colframe=black,
    arc=0mm,
    title=Pseudo-code
]
\tiny
\begin{lstlisting}
HEAP-EXTRACT-MAX(A, n)

    if n < 1

        error "heap underflow"


    max = A[1]

    A[1] = A[n]

    n = n - 1


    MAX-HEAPIFY(A, 1, n)


    return max
\end{lstlisting}
\end{tcolorbox}


\begin{tcolorbox}[
    colback=cyan!20,
    colframe=black,
    arc=0mm
]
\small
\textbf{Time Complexity:}
$O(\log n)$
\end{tcolorbox}


\begin{tcolorbox}[
    colback=violet!20,
    colframe=black,
    arc=0mm
]
\small
\textbf{Space Complexity:}
$O(1)$ auxiliary space.
\end{tcolorbox}


\begin{tcolorbox}[
    colback=pink!20,
    colframe=black,
    arc=0mm
]
\small
\textbf{Recurrence Formula:}
$T(n)=T(n/2)+\Theta(1)$
$\Rightarrow O(\log n)$
\end{tcolorbox}

\end{tcolorbox}

















\begin{tcolorbox}[
    colback=gray!30,
    colframe=black,
    boxrule=1pt,
    arc=0mm
]

\begin{tcolorbox}[
    colback=gray!10,
    colframe=black,
    arc=0mm
]
\centering
{\Huge\bfseries Heap Maximum}
\end{tcolorbox}


\begin{tcolorbox}[
    colback=gray!10,
    colframe=black,
    arc=0mm
]
\small
Return the maximum element stored in a max-heap.
Since the heap property guarantees that the largest element is always at the root,
the operation simply accesses the first element of the array.
\end{tcolorbox}


\begin{tcolorbox}[
    colback=black!5,
    colframe=black,
    arc=0mm,
    title=Pseudo-code
]
\tiny
\begin{lstlisting}
HEAP-MAXIMUM(A)

    return A[1]
\end{lstlisting}
\end{tcolorbox}


\begin{tcolorbox}[
    colback=cyan!20,
    colframe=black,
    arc=0mm
]
\small
\textbf{Time Complexity:}
$O(1)$
\end{tcolorbox}


\begin{tcolorbox}[
    colback=violet!20,
    colframe=black,
    arc=0mm
]
\small
\textbf{Space Complexity:}
$O(1)$
\end{tcolorbox}


\begin{tcolorbox}[
    colback=pink!20,
    colframe=black,
    arc=0mm
]
\small
\textbf{Recurrence Formula:}
No recurrence formula (constant-time operation).
\end{tcolorbox}

\end{tcolorbox}





















\begin{tcolorbox}[
    colback=gray!30,
    colframe=black,
    boxrule=1pt,
    arc=0mm
]

\begin{tcolorbox}[
    colback=gray!10,
    colframe=black,
    arc=0mm
]
\centering
{\Huge\bfseries Heap Increase Key}
\end{tcolorbox}


\begin{tcolorbox}[
    colback=gray!10,
    colframe=black,
    arc=0mm
]
\small
Update the value of a node in a max-heap and restore the heap property.
After increasing the key, the element is repeatedly swapped with its parent
until it reaches a valid position.
\end{tcolorbox}


\begin{tcolorbox}[
    colback=black!5,
    colframe=black,
    arc=0mm,
    title=Pseudo-code
]
\tiny
\begin{lstlisting}
HEAP-INCREASE-KEY(A, i, key)

    if key < A[i]

        error "new key is smaller than current key"


    A[i] = key


    while i > 1 and A[PARENT(i)] < A[i]

        exchange A[i] and A[PARENT(i)]

        i = PARENT(i)
\end{lstlisting}
\end{tcolorbox}


\begin{tcolorbox}[
    colback=cyan!20,
    colframe=black,
    arc=0mm
]
\small
\textbf{Time Complexity:}
$O(\log n)$
\end{tcolorbox}


\begin{tcolorbox}[
    colback=violet!20,
    colframe=black,
    arc=0mm
]
\small
\textbf{Space Complexity:}
$O(1)$
\end{tcolorbox}


\begin{tcolorbox}[
    colback=pink!20,
    colframe=black,
    arc=0mm
]
\small
\textbf{Recurrence Formula:}
No recurrence formula (iterative upward traversal).
\end{tcolorbox}

\end{tcolorbox}
































\begin{tcolorbox}[
    colback=gray!30,
    colframe=black,
    boxrule=1pt,
    arc=0mm
]

\begin{tcolorbox}[
    colback=gray!10,
    colframe=black,
    arc=0mm
]
\centering
{\Huge\bfseries Max Heap Insert}
\end{tcolorbox}


\begin{tcolorbox}[
    colback=gray!10,
    colframe=black,
    arc=0mm
]
\small
Insert a new element into a max-heap while preserving the heap property.
A temporary value of $8$ is first inserted, then HEAP-INCREASE-KEY
moves the new key upward until it reaches its correct position.
\end{tcolorbox}


\begin{tcolorbox}[
    colback=black!5,
    colframe=black,
    arc=0mm,
    title=Pseudo-code
]
\tiny
\begin{lstlisting}
MAX-HEAP-INSERT(A, key, n)

    n = n + 1

    A[n] = 8

    HEAP-INCREASE-KEY(A, n, key)
\end{lstlisting}
\end{tcolorbox}


\begin{tcolorbox}[
    colback=cyan!20,
    colframe=black,
    arc=0mm
]
\small
\textbf{Time Complexity:}
$O(\log n)$
\end{tcolorbox}


\begin{tcolorbox}[
    colback=violet!20,
    colframe=black,
    arc=0mm
]
\small
\textbf{Space Complexity:}
$O(1)$
\end{tcolorbox}


\begin{tcolorbox}[
    colback=pink!20,
    colframe=black,
    arc=0mm
]
\small
\textbf{Recurrence Formula:}
No recurrence formula (uses HEAP-INCREASE-KEY).
\end{tcolorbox}

\end{tcolorbox}

























\begin{tcolorbox}[
    colback=gray!30,
    colframe=black,
    boxrule=1pt,
    arc=0mm
]

\begin{tcolorbox}[
    colback=gray!10,
    colframe=black,
    arc=0mm
]
\centering
{\Huge\bfseries Stack Empty}
\end{tcolorbox}


\begin{tcolorbox}[
    colback=gray!10,
    colframe=black,
    arc=0mm
]
\small
Check whether a stack contains any elements.
The operation simply tests the value of the top pointer.
\end{tcolorbox}


\begin{tcolorbox}[
    colback=black!5,
    colframe=black,
    arc=0mm,
    title=Pseudo-code
]
\tiny
\begin{lstlisting}
STACK-EMPTY(S)

    if S.top == 0

        return TRUE

    else

        return FALSE
\end{lstlisting}
\end{tcolorbox}


\begin{tcolorbox}[
    colback=cyan!20,
    colframe=black,
    arc=0mm
]
\small
\textbf{Time Complexity:}
$O(1)$
\end{tcolorbox}


\begin{tcolorbox}[
    colback=violet!20,
    colframe=black,
    arc=0mm
]
\small
\textbf{Space Complexity:}
$O(1)$
\end{tcolorbox}


\begin{tcolorbox}[
    colback=pink!20,
    colframe=black,
    arc=0mm
]
\small
\textbf{Recurrence Formula:}
No recurrence formula (constant-time operation).
\end{tcolorbox}

\end{tcolorbox}


























\begin{tcolorbox}[
    colback=gray!30,
    colframe=black,
    boxrule=1pt,
    arc=0mm
]

\begin{tcolorbox}[
    colback=gray!10,
    colframe=black,
    arc=0mm
]
\centering
{\Huge\bfseries Stack Push}
\end{tcolorbox}


\begin{tcolorbox}[
    colback=gray!10,
    colframe=black,
    arc=0mm
]
\small
Insert a new element at the top of the stack.
The operation increments the top pointer and stores the value.
\end{tcolorbox}


\begin{tcolorbox}[
    colback=black!5,
    colframe=black,
    arc=0mm,
    title=Pseudo-code
]
\tiny
\begin{lstlisting}
PUSH(S, x)

    S.top = S.top + 1

    S[S.top] = x
\end{lstlisting}
\end{tcolorbox}


\begin{tcolorbox}[
    colback=cyan!20,
    colframe=black,
    arc=0mm
]
\small
\textbf{Time Complexity:}
$O(1)$
\end{tcolorbox}


\begin{tcolorbox}[
    colback=violet!20,
    colframe=black,
    arc=0mm
]
\small
\textbf{Space Complexity:}
$O(1)$
\end{tcolorbox}


\begin{tcolorbox}[
    colback=pink!20,
    colframe=black,
    arc=0mm
]
\small
\textbf{Recurrence Formula:}
No recurrence formula (constant-time operation).
\end{tcolorbox}

\end{tcolorbox}


















\begin{tcolorbox}[
    colback=gray!30,
    colframe=black,
    boxrule=1pt,
    arc=0mm
]

\begin{tcolorbox}[
    colback=gray!10,
    colframe=black,
    arc=0mm
]
\centering
{\Huge\bfseries Stack Pop}
\end{tcolorbox}


\begin{tcolorbox}[
    colback=gray!10,
    colframe=black,
    arc=0mm
]
\small
Remove and return the top element of the stack.
The operation checks if the stack is empty, then decreases the top pointer.
\end{tcolorbox}


\begin{tcolorbox}[
    colback=black!5,
    colframe=black,
    arc=0mm,
    title=Pseudo-code
]
\tiny
\begin{lstlisting}
POP(S)

    if STACK-EMPTY(S)

        error "underflow"

    else

        S.top = S.top - 1

        return S[S.top + 1]
\end{lstlisting}
\end{tcolorbox}


\begin{tcolorbox}[
    colback=cyan!20,
    colframe=black,
    arc=0mm
]
\small
\textbf{Time Complexity:}
$O(1)$
\end{tcolorbox}


\begin{tcolorbox}[
    colback=violet!20,
    colframe=black,
    arc=0mm
]
\small
\textbf{Space Complexity:}
$O(1)$
\end{tcolorbox}


\begin{tcolorbox}[
    colback=pink!20,
    colframe=black,
    arc=0mm
]
\small
\textbf{Recurrence Formula:}
No recurrence formula (constant-time operation).
\end{tcolorbox}

\end{tcolorbox}



















\begin{tcolorbox}[
    colback=gray!30,
    colframe=black,
    boxrule=1pt,
    arc=0mm
]

\begin{tcolorbox}[
    colback=gray!10,
    colframe=black,
    arc=0mm
]
\centering
{\Huge\bfseries Enqueue}
\end{tcolorbox}


\begin{tcolorbox}[
    colback=gray!10,
    colframe=black,
    arc=0mm
]
\small
Insert a new element at the rear of a queue.
The element is stored at the tail position, then the tail pointer
is updated using circular array indexing.
\end{tcolorbox}


\begin{tcolorbox}[
    colback=black!5,
    colframe=black,
    arc=0mm,
    title=Pseudo-code
]
\tiny
\begin{lstlisting}
ENQUEUE(Q, x)

    Q[Q.tail] = x


    if Q.tail == Q.length

        Q.tail = 1

    else

        Q.tail = Q.tail + 1
\end{lstlisting}
\end{tcolorbox}


\begin{tcolorbox}[
    colback=cyan!20,
    colframe=black,
    arc=0mm
]
\small
\textbf{Time Complexity:}
$O(1)$
\end{tcolorbox}


\begin{tcolorbox}[
    colback=violet!20,
    colframe=black,
    arc=0mm
]
\small
\textbf{Space Complexity:}
$O(1)$
\end{tcolorbox}


\begin{tcolorbox}[
    colback=pink!20,
    colframe=black,
    arc=0mm
]
\small
\textbf{Recurrence Formula:}
No recurrence formula (constant-time operation).
\end{tcolorbox}

\end{tcolorbox}



























\begin{tcolorbox}[
    colback=gray!30,
    colframe=black,
    boxrule=1pt,
    arc=0mm
]

\begin{tcolorbox}[
    colback=gray!10,
    colframe=black,
    arc=0mm
]
\centering
{\Huge\bfseries Dequeue}
\end{tcolorbox}


\begin{tcolorbox}[
    colback=gray!10,
    colframe=black,
    arc=0mm
]
\small
Remove and return the element at the front of a queue.
The element at the head is read, then the head pointer is advanced
using circular array indexing.
\end{tcolorbox}


\begin{tcolorbox}[
    colback=black!5,
    colframe=black,
    arc=0mm,
    title=Pseudo-code
]
\tiny
\begin{lstlisting}
DEQUEUE(Q)

    x = Q[Q.head]


    if Q.head == Q.length

        Q.head = 1

    else

        Q.head = Q.head + 1


    return x
\end{lstlisting}
\end{tcolorbox}


\begin{tcolorbox}[
    colback=cyan!20,
    colframe=black,
    arc=0mm
]
\small
\textbf{Time Complexity:}
$O(1)$
\end{tcolorbox}


\begin{tcolorbox}[
    colback=violet!20,
    colframe=black,
    arc=0mm
]
\small
\textbf{Space Complexity:}
$O(1)$
\end{tcolorbox}


\begin{tcolorbox}[
    colback=pink!20,
    colframe=black,
    arc=0mm
]
\small
\textbf{Recurrence Formula:}
No recurrence formula (constant-time operation).
\end{tcolorbox}

\end{tcolorbox}




























\begin{tcolorbox}[
    colback=gray!30,
    colframe=black,
    boxrule=1pt,
    arc=0mm
]

\begin{tcolorbox}[
    colback=gray!10,
    colframe=black,
    arc=0mm
]
\centering
{\Huge\bfseries Linked List Search}
\end{tcolorbox}


\begin{tcolorbox}[
    colback=gray!10,
    colframe=black,
    arc=0mm
]
\small
Search for an element with a given key in a linked list.
The algorithm starts from the head and traverses each node until
the key is found or the end of the list is reached.
\end{tcolorbox}


\begin{tcolorbox}[
    colback=black!5,
    colframe=black,
    arc=0mm,
    title=Pseudo-code
]
\tiny
\begin{lstlisting}
LIST-SEARCH(L, k)

    x = L.head


    while x != NIL and x.key != k

        x = x.next


    return x
\end{lstlisting}
\end{tcolorbox}


\begin{tcolorbox}[
    colback=cyan!20,
    colframe=black,
    arc=0mm
]
\small
\textbf{Time Complexity:}
$O(n)$
\end{tcolorbox}


\begin{tcolorbox}[
    colback=violet!20,
    colframe=black,
    arc=0mm
]
\small
\textbf{Space Complexity:}
$O(1)$
\end{tcolorbox}


\begin{tcolorbox}[
    colback=pink!20,
    colframe=black,
    arc=0mm
]
\small
\textbf{Recurrence Formula:}
No recurrence formula (iterative traversal).
\end{tcolorbox}

\end{tcolorbox}




















\begin{tcolorbox}[
    colback=gray!30,
    colframe=black,
    boxrule=1pt,
    arc=0mm
]

\begin{tcolorbox}[
    colback=gray!10,
    colframe=black,
    arc=0mm
]
\centering
{\Huge\bfseries Linked List Insert}
\end{tcolorbox}


\begin{tcolorbox}[
    colback=gray!10,
    colframe=black,
    arc=0mm
]
\small
Insert a new node at the beginning of a doubly linked list.
The operation updates the pointers so that the new node becomes
the new head while preserving the list connections.
\end{tcolorbox}


\begin{tcolorbox}[
    colback=black!5,
    colframe=black,
    arc=0mm,
    title=Pseudo-code
]
\tiny
\begin{lstlisting}
LIST-INSERT(L, x)

    x.next = L.head


    if L.head != NIL

        L.head.prev = x


    L.head = x

    x.prev = NIL
\end{lstlisting}
\end{tcolorbox}


\begin{tcolorbox}[
    colback=cyan!20,
    colframe=black,
    arc=0mm
]
\small
\textbf{Time Complexity:}
$O(1)$
\end{tcolorbox}


\begin{tcolorbox}[
    colback=violet!20,
    colframe=black,
    arc=0mm
]
\small
\textbf{Space Complexity:}
$O(1)$
\end{tcolorbox}


\begin{tcolorbox}[
    colback=pink!20,
    colframe=black,
    arc=0mm
]
\small
\textbf{Recurrence Formula:}
No recurrence formula (constant-time pointer update).
\end{tcolorbox}

\end{tcolorbox}






























\begin{tcolorbox}[
    colback=gray!30,
    colframe=black,
    boxrule=1pt,
    arc=0mm
]

\begin{tcolorbox}[
    colback=gray!10,
    colframe=black,
    arc=0mm
]
\centering
{\Huge\bfseries Linked List Delete}
\end{tcolorbox}


\begin{tcolorbox}[
    colback=gray!10,
    colframe=black,
    arc=0mm
]
\small
Remove a node from a doubly linked list.
The algorithm updates the neighboring nodes' pointers so the list
remains correctly connected after the deletion.
\end{tcolorbox}


\begin{tcolorbox}[
    colback=black!5,
    colframe=black,
    arc=0mm,
    title=Pseudo-code
]
\tiny
\begin{lstlisting}
LIST-DELETE(L, x)

    if x.prev != NIL

        x.prev.next = x.next

    else

        L.head = x.next



    if x.next != NIL

        x.next.prev = x.prev
\end{lstlisting}
\end{tcolorbox}


\begin{tcolorbox}[
    colback=cyan!20,
    colframe=black,
    arc=0mm
]
\small
\textbf{Time Complexity:}
$O(1)$
\end{tcolorbox}


\begin{tcolorbox}[
    colback=violet!20,
    colframe=black,
    arc=0mm
]
\small
\textbf{Space Complexity:}
$O(1)$
\end{tcolorbox}


\begin{tcolorbox}[
    colback=pink!20,
    colframe=black,
    arc=0mm
]
\small
\textbf{Recurrence Formula:}
No recurrence formula (constant-time pointer update).
\end{tcolorbox}

\end{tcolorbox}






















\begin{tcolorbox}[
    colback=gray!30,
    colframe=black,
    boxrule=1pt,
    arc=0mm
]

\begin{tcolorbox}[
    colback=gray!10,
    colframe=black,
    arc=0mm
]
\centering
{\Huge\bfseries Tree Search}
\end{tcolorbox}


\begin{tcolorbox}[
    colback=gray!10,
    colframe=black,
    arc=0mm
]
\small
Search for a node with a given key in a binary search tree.
The algorithm follows the BST property: go left if the key is smaller,
otherwise go right, until the key is found or the tree ends.
\end{tcolorbox}


\begin{tcolorbox}[
    colback=black!5,
    colframe=black,
    arc=0mm,
    title=Pseudo-code
]
\tiny
\begin{lstlisting}
TREE-SEARCH(x, k)

    if x == NIL or k == x.key

        return x


    if k < x.key

        return TREE-SEARCH(x.left, k)

    else

        return TREE-SEARCH(x.right, k)
\end{lstlisting}
\end{tcolorbox}


\begin{tcolorbox}[
    colback=cyan!20,
    colframe=black,
    arc=0mm
]
\small
\textbf{Time Complexity:}
$O(h)$ where $h$ is the tree height.
\end{tcolorbox}


\begin{tcolorbox}[
    colback=violet!20,
    colframe=black,
    arc=0mm
]
\small
\textbf{Space Complexity:}
$O(h)$ due to recursive call stack.
\end{tcolorbox}


\begin{tcolorbox}[
    colback=pink!20,
    colframe=black,
    arc=0mm
]
\small
\textbf{Recurrence Formula:}
$T(h)=T(h-1)+\Theta(1)$
$\Rightarrow O(h)$
\end{tcolorbox}

\end{tcolorbox}


















\begin{tcolorbox}[
    colback=gray!30,
    colframe=black,
    boxrule=1pt,
    arc=0mm
]

\begin{tcolorbox}[
    colback=gray!10,
    colframe=black,
    arc=0mm
]
\centering
{\Huge\bfseries Tree Minimum}
\end{tcolorbox}


\begin{tcolorbox}[
    colback=gray!10,
    colframe=black,
    arc=0mm
]
\small
Find the smallest element in a binary search tree.
The minimum value is always located by following left children
until reaching a node without a left child.
\end{tcolorbox}


\begin{tcolorbox}[
    colback=black!5,
    colframe=black,
    arc=0mm,
    title=Pseudo-code
]
\tiny
\begin{lstlisting}
TREE-MINIMUM(x)

    while x.left != NIL

        x = x.left


    return x
\end{lstlisting}
\end{tcolorbox}


\begin{tcolorbox}[
    colback=cyan!20,
    colframe=black,
    arc=0mm
]
\small
\textbf{Time Complexity:}
$O(h)$ where $h$ is the tree height.
\end{tcolorbox}


\begin{tcolorbox}[
    colback=violet!20,
    colframe=black,
    arc=0mm
]
\small
\textbf{Space Complexity:}
$O(1)$
\end{tcolorbox}


\begin{tcolorbox}[
    colback=pink!20,
    colframe=black,
    arc=0mm
]
\small
\textbf{Recurrence Formula:}
No recurrence formula (iterative traversal).
\end{tcolorbox}

\end{tcolorbox}



























\begin{tcolorbox}[
    colback=gray!30,
    colframe=black,
    boxrule=1pt,
    arc=0mm
]

\begin{tcolorbox}[
    colback=gray!10,
    colframe=black,
    arc=0mm
]
\centering
{\Huge\bfseries Tree Maximum}
\end{tcolorbox}


\begin{tcolorbox}[
    colback=gray!10,
    colframe=black,
    arc=0mm
]
\small
Find the largest element in a binary search tree.
The maximum value is located by following right children
until reaching a node without a right child.
\end{tcolorbox}


\begin{tcolorbox}[
    colback=black!5,
    colframe=black,
    arc=0mm,
    title=Pseudo-code
]
\tiny
\begin{lstlisting}
TREE-MAXIMUM(x)

    while x.right != NIL

        x = x.right


    return x
\end{lstlisting}
\end{tcolorbox}


\begin{tcolorbox}[
    colback=cyan!20,
    colframe=black,
    arc=0mm
]
\small
\textbf{Time Complexity:}
$O(h)$ where $h$ is the tree height.
\end{tcolorbox}


\begin{tcolorbox}[
    colback=violet!20,
    colframe=black,
    arc=0mm
]
\small
\textbf{Space Complexity:}
$O(1)$
\end{tcolorbox}


\begin{tcolorbox}[
    colback=pink!20,
    colframe=black,
    arc=0mm
]
\small
\textbf{Recurrence Formula:}
No recurrence formula (iterative traversal).
\end{tcolorbox}

\end{tcolorbox}




























\begin{tcolorbox}[
    colback=gray!30,
    colframe=black,
    boxrule=1pt,
    arc=0mm
]

\begin{tcolorbox}[
    colback=gray!10,
    colframe=black,
    arc=0mm
]
\centering
{\Huge\bfseries Tree Successor}
\end{tcolorbox}


\begin{tcolorbox}[
    colback=gray!10,
    colframe=black,
    arc=0mm
]
\small
Find the next node in the in-order traversal of a binary search tree.
If a right subtree exists, the successor is the minimum of that subtree.
Otherwise, move upward until finding a node that is a left child of its parent.
\end{tcolorbox}


\begin{tcolorbox}[
    colback=black!5,
    colframe=black,
    arc=0mm,
    title=Pseudo-code
]
\tiny
\begin{lstlisting}
TREE-SUCCESSOR(x)

    if x.right != NIL

        return TREE-MINIMUM(x.right)


    y = x.p


    while y != NIL and x == y.right

        x = y

        y = y.p


    return y
\end{lstlisting}
\end{tcolorbox}


\begin{tcolorbox}[
    colback=cyan!20,
    colframe=black,
    arc=0mm
]
\small
\textbf{Time Complexity:}
$O(h)$ where $h$ is the tree height.
\end{tcolorbox}


\begin{tcolorbox}[
    colback=violet!20,
    colframe=black,
    arc=0mm
]
\small
\textbf{Space Complexity:}
$O(1)$
\end{tcolorbox}


\begin{tcolorbox}[
    colback=pink!20,
    colframe=black,
    arc=0mm
]
\small
\textbf{Recurrence Formula:}
No recurrence formula (iterative tree traversal).
\end{tcolorbox}

\end{tcolorbox}
























\begin{tcolorbox}[
    colback=gray!30,
    colframe=black,
    boxrule=1pt,
    arc=0mm
]

\begin{tcolorbox}[
    colback=gray!10,
    colframe=black,
    arc=0mm
]
\centering
{\Huge\bfseries Inorder Tree Walk}
\end{tcolorbox}


\begin{tcolorbox}[
    colback=gray!10,
    colframe=black,
    arc=0mm
]
\small
Traverse all nodes of a binary search tree in sorted order.
The algorithm follows the in-order strategy:
left subtree, node, then right subtree.
\end{tcolorbox}


\begin{tcolorbox}[
    colback=black!5,
    colframe=black,
    arc=0mm,
    title=Pseudo-code
]
\tiny
\begin{lstlisting}
INORDER-TREE-WALK(x)

    if x != NIL

        INORDER-TREE-WALK(x.left)

        print key[x]

        INORDER-TREE-WALK(x.right)
\end{lstlisting}
\end{tcolorbox}


\begin{tcolorbox}[
    colback=cyan!20,
    colframe=black,
    arc=0mm
]
\small
\textbf{Time Complexity:}
$O(n)$
\end{tcolorbox}


\begin{tcolorbox}[
    colback=violet!20,
    colframe=black,
    arc=0mm
]
\small
\textbf{Space Complexity:}
$O(h)$ recursion stack.
\end{tcolorbox}


\begin{tcolorbox}[
    colback=pink!20,
    colframe=black,
    arc=0mm
]
\small
\textbf{Recurrence Formula:}
$T(n)=T(n_L)+T(n_R)+\Theta(1)$
$\Rightarrow O(n)$
\end{tcolorbox}

\end{tcolorbox}

























\begin{tcolorbox}[
    colback=gray!30,
    colframe=black,
    boxrule=1pt,
    arc=0mm
]

\begin{tcolorbox}[
    colback=gray!10,
    colframe=black,
    arc=0mm
]
\centering
{\Huge\bfseries Preorder Tree Walk}
\end{tcolorbox}


\begin{tcolorbox}[
    colback=gray!10,
    colframe=black,
    arc=0mm
]
\small
Traverse all nodes of a binary tree by visiting the root first,
then the left subtree, and finally the right subtree.
This allows processing a node before its children.
\end{tcolorbox}


\begin{tcolorbox}[
    colback=black!5,
    colframe=black,
    arc=0mm,
    title=Pseudo-code
]
\tiny
\begin{lstlisting}
PREORDER-TREE-WALK(x)

    if x != NIL

        print key[x]

        PREORDER-TREE-WALK(x.left)

        PREORDER-TREE-WALK(x.right)
\end{lstlisting}
\end{tcolorbox}


\begin{tcolorbox}[
    colback=cyan!20,
    colframe=black,
    arc=0mm
]
\small
\textbf{Time Complexity:}
$O(n)$
\end{tcolorbox}


\begin{tcolorbox}[
    colback=violet!20,
    colframe=black,
    arc=0mm
]
\small
\textbf{Space Complexity:}
$O(h)$ recursion stack.
\end{tcolorbox}


\begin{tcolorbox}[
    colback=pink!20,
    colframe=black,
    arc=0mm
]
\small
\textbf{Recurrence Formula:}
$T(n)=T(n_L)+T(n_R)+\Theta(1)$
$\Rightarrow O(n)$
\end{tcolorbox}

\end{tcolorbox}


























\begin{tcolorbox}[
    colback=gray!30,
    colframe=black,
    boxrule=1pt,
    arc=0mm
]

\begin{tcolorbox}[
    colback=gray!10,
    colframe=black,
    arc=0mm
]
\centering
{\Huge\bfseries Postorder Tree Walk}
\end{tcolorbox}


\begin{tcolorbox}[
    colback=gray!10,
    colframe=black,
    arc=0mm
]
\small
Traverse all nodes of a binary tree by visiting the children first
and the root last.
The traversal order is: left subtree, right subtree, then node.
\end{tcolorbox}


\begin{tcolorbox}[
    colback=black!5,
    colframe=black,
    arc=0mm,
    title=Pseudo-code
]
\tiny
\begin{lstlisting}
POSTORDER-TREE-WALK(x)

    if x != NIL

        POSTORDER-TREE-WALK(x.left)

        POSTORDER-TREE-WALK(x.right)

        print key[x]
\end{lstlisting}
\end{tcolorbox}


\begin{tcolorbox}[
    colback=cyan!20,
    colframe=black,
    arc=0mm
]
\small
\textbf{Time Complexity:}
$O(n)$
\end{tcolorbox}


\begin{tcolorbox}[
    colback=violet!20,
    colframe=black,
    arc=0mm
]
\small
\textbf{Space Complexity:}
$O(h)$ recursion stack.
\end{tcolorbox}


\begin{tcolorbox}[
    colback=pink!20,
    colframe=black,
    arc=0mm
]
\small
\textbf{Recurrence Formula:}
$T(n)=T(n_L)+T(n_R)+\Theta(1)$
$\Rightarrow O(n)$
\end{tcolorbox}

\end{tcolorbox}































\begin{tcolorbox}[
    colback=gray!30,
    colframe=black,
    boxrule=1pt,
    arc=0mm
]

\begin{tcolorbox}[
    colback=gray!10,
    colframe=black,
    arc=0mm
]
\centering
{\Huge\bfseries Tree Insert}
\end{tcolorbox}


\begin{tcolorbox}[
    colback=gray!10,
    colframe=black,
    arc=0mm
]
\small
Insert a new node into a binary search tree while preserving the BST property.
The algorithm searches for the correct position and attaches the node as a leaf.
\end{tcolorbox}


\begin{tcolorbox}[
    colback=black!5,
    colframe=black,
    arc=0mm,
    title=Pseudo-code
]
\tiny
\begin{lstlisting}
TREE-INSERT(T, z)

    y = NIL

    x = T.root


    while x != NIL

        y = x

        if z.key < x.key

            x = x.left

        else

            x = x.right


    z.p = y


    if y == NIL

        T.root = z

    elseif z.key < y.key

        y.left = z

    else

        y.right = z
\end{lstlisting}
\end{tcolorbox}


\begin{tcolorbox}[
    colback=cyan!20,
    colframe=black,
    arc=0mm
]
\small
\textbf{Time Complexity:}
$O(h)$ where $h$ is the tree height.
\end{tcolorbox}


\begin{tcolorbox}[
    colback=violet!20,
    colframe=black,
    arc=0mm
]
\small
\textbf{Space Complexity:}
$O(1)$
\end{tcolorbox}


\begin{tcolorbox}[
    colback=pink!20,
    colframe=black,
    arc=0mm
]
\small
\textbf{Recurrence Formula:}
No recurrence formula (iterative tree traversal).
\end{tcolorbox}

\end{tcolorbox}



























\begin{tcolorbox}[
    colback=gray!30,
    colframe=black,
    boxrule=1pt,
    arc=0mm
]

\begin{tcolorbox}[
    colback=gray!10,
    colframe=black,
    arc=0mm
]
\centering
{\Huge\bfseries Tree Delete}
\end{tcolorbox}


\begin{tcolorbox}[
    colback=gray!10,
    colframe=black,
    arc=0mm
]
\small
Remove a node from a binary search tree while preserving the BST property.
The algorithm handles the three cases: no child, one child, or two children.
It uses TRANSPLANT to safely reconnect subtrees.
\end{tcolorbox}


\begin{tcolorbox}[
    colback=black!5,
    colframe=black,
    arc=0mm,
    title=Pseudo-code
]
\tiny
\begin{lstlisting}
TRANSPLANT(T, u, v)

    if u.p == NIL

        T.root = v

    elseif u == u.p.left

        u.p.left = v

    else

        u.p.right = v


    if v != NIL

        v.p = u.p



TREE-DELETE(T, z)

    if z.left == NIL

        TRANSPLANT(T, z, z.right)


    elseif z.right == NIL

        TRANSPLANT(T, z, z.left)


    else

        y = TREE-MINIMUM(z.right)


        if y.p != z

            TRANSPLANT(T, y, y.right)

            y.right = z.right

            y.right.p = y


        TRANSPLANT(T, z, y)

        y.left = z.left

        y.left.p = y
\end{lstlisting}
\end{tcolorbox}


\begin{tcolorbox}[
    colback=cyan!20,
    colframe=black,
    arc=0mm
]
\small
\textbf{Time Complexity:}
$O(h)$ where $h$ is the tree height.
\end{tcolorbox}


\begin{tcolorbox}[
    colback=violet!20,
    colframe=black,
    arc=0mm
]
\small
\textbf{Space Complexity:}
$O(1)$
\end{tcolorbox}


\begin{tcolorbox}[
    colback=pink!20,
    colframe=black,
    arc=0mm
]
\small
\textbf{Recurrence Formula:}
No recurrence formula (iterative tree modification).
\end{tcolorbox}

\end{tcolorbox}



























\begin{tcolorbox}[
    colback=gray!30,
    colframe=black,
    boxrule=1pt,
    arc=0mm
]

\begin{tcolorbox}[
    colback=gray!10,
    colframe=black,
    arc=0mm
]
\centering
{\Huge\bfseries Rod Cutting Algorithm}
\end{tcolorbox}


\begin{tcolorbox}[
    colback=gray!10,
    colframe=black,
    arc=0mm
]
\small
Find the maximum revenue obtainable by cutting a rod of length $n$.
The dynamic programming approach uses memoization to avoid recomputing
the same subproblems.
\end{tcolorbox}


\begin{tcolorbox}[
    colback=black!5,
    colframe=black,
    arc=0mm,
    title=Pseudo-code
]
\tiny
\begin{lstlisting}
Memoized-Cut-Rod(p, n)

    let r[0..n] be a new array

    for i = 0 to n

        r[i] = 8


    return Memoized-Cut-Rod-Aux(p, n, r)
\end{lstlisting}
\end{tcolorbox}


\begin{tcolorbox}[
    colback=cyan!20,
    colframe=black,
    arc=0mm
]
\small
\textbf{Time Complexity:}
$O(n^2)$
\end{tcolorbox}


\begin{tcolorbox}[
    colback=violet!20,
    colframe=black,
    arc=0mm
]
\small
\textbf{Space Complexity:}
$O(n)$
\end{tcolorbox}


\begin{tcolorbox}[
    colback=pink!20,
    colframe=black,
    arc=0mm
]
\small
\textbf{Recurrence Formula:}
$T(n)=\sum_{j=1}^{n}T(n-j)+O(n)$
$\Rightarrow O(n^2)$ with memoization.
\end{tcolorbox}

\end{tcolorbox}





























\begin{tcolorbox}[
    colback=gray!30,
    colframe=black,
    boxrule=1pt,
    arc=0mm
]

\begin{tcolorbox}[
    colback=gray!10,
    colframe=black,
    arc=0mm
]
\centering
{\Huge\bfseries Memoized Cut Rod Aux}
\end{tcolorbox}


\begin{tcolorbox}[
    colback=gray!10,
    colframe=black,
    arc=0mm
]
\small
Compute the optimal revenue for a rod of length $n$ using memoization.
The algorithm stores previously computed solutions to avoid solving the same
subproblems multiple times.
\end{tcolorbox}


\begin{tcolorbox}[
    colback=black!5,
    colframe=black,
    arc=0mm,
    title=Pseudo-code
]
\tiny
\begin{lstlisting}
Memoized-Cut-Rod-Aux(p, n, r)

    if r[n] >= 0

        return r[n]


    if n == 0

        q = 0

    else

        q = 8

        for i = 1 to n

            q = max(q,
                p[i] +
                Memoized-Cut-Rod-Aux(p, n-i, r))


    r[n] = q

    return q
\end{lstlisting}
\end{tcolorbox}


\begin{tcolorbox}[
    colback=cyan!20,
    colframe=black,
    arc=0mm
]
\small
\textbf{Time Complexity:}
$O(n^2)$
\end{tcolorbox}


\begin{tcolorbox}[
    colback=violet!20,
    colframe=black,
    arc=0mm
]
\small
\textbf{Space Complexity:}
$O(n)$
\end{tcolorbox}


\begin{tcolorbox}[
    colback=pink!20,
    colframe=black,
    arc=0mm
]
\small
\textbf{Recurrence Formula:}
$T(n)=\sum_{i=1}^{n}T(n-i)+O(n)$
with memoization $\Rightarrow O(n^2)$.
\end{tcolorbox}

\end{tcolorbox}





















\begin{tcolorbox}[
    colback=gray!30,
    colframe=black,
    boxrule=1pt,
    arc=0mm
]

\begin{tcolorbox}[
    colback=gray!10,
    colframe=black,
    arc=0mm
]
\centering
{\Huge\bfseries Bottom-Up Cut Rod}
\end{tcolorbox}


\begin{tcolorbox}[
    colback=gray!10,
    colframe=black,
    arc=0mm
]
\small
Compute the maximum revenue obtainable from a rod of length $n$
using bottom-up dynamic programming.
The algorithm solves smaller subproblems first and stores their results
to build the final optimal solution.
\end{tcolorbox}


\begin{tcolorbox}[
    colback=black!5,
    colframe=black,
    arc=0mm,
    title=Pseudo-code
]
\tiny
\begin{lstlisting}
BOTTOM-UP-CUT-ROD(p, n)

    let r[0..n] be a new array

    r[0] = 0


    for j = 1 to n

        q = 8

        for i = 1 to j

            q = max(q,
                    p[i] + r[j-i])


        r[j] = q


    return r[n]
\end{lstlisting}
\end{tcolorbox}


\begin{tcolorbox}[
    colback=cyan!20,
    colframe=black,
    arc=0mm
]
\small
\textbf{Time Complexity:}
$O(n^2)$
\end{tcolorbox}


\begin{tcolorbox}[
    colback=violet!20,
    colframe=black,
    arc=0mm
]
\small
\textbf{Space Complexity:}
$O(n)$
\end{tcolorbox}


\begin{tcolorbox}[
    colback=pink!20,
    colframe=black,
    arc=0mm
]
\small
\textbf{Recurrence Formula:}
$T(n)=\sum_{i=1}^{n}T(n-i)+\Theta(n)$
$\Rightarrow O(n^2)$ with dynamic programming.
\end{tcolorbox}

\end{tcolorbox}























\begin{tcolorbox}[
    colback=gray!30,
    colframe=black,
    boxrule=1pt,
    arc=0mm
]

\begin{tcolorbox}[
    colback=gray!10,
    colframe=black,
    arc=0mm
]
\centering
{\Huge\bfseries Matrix Chain Order}
\end{tcolorbox}


\begin{tcolorbox}[
    colback=gray!10,
    colframe=black,
    arc=0mm
]
\small
Find the optimal parenthesization of a matrix chain to minimize
the number of scalar multiplications.
The algorithm uses dynamic programming by storing optimal costs
for smaller subchains.
\end{tcolorbox}


\begin{tcolorbox}[
    colback=black!5,
    colframe=black,
    arc=0mm,
    title=Pseudo-code
]
\tiny
\begin{lstlisting}
MATRIX-CHAIN-ORDER(p)

    n = p.length - 1

    let m[1..n,1..n] and s[1..n,1..n]


    for i = 1 to n

        m[i,i] = 0


    for l = 2 to n

        for i = 1 to n-l+1

            j = i+l-1

            m[i,j] = 8


            for k = i to j-1

                q = m[i,k]
                    + m[k+1,j]
                    + p[i-1]*p[k]*p[j]


                if q < m[i,j]

                    m[i,j] = q

                    s[i,j] = k


    return m and s
\end{lstlisting}
\end{tcolorbox}


\begin{tcolorbox}[
    colback=cyan!20,
    colframe=black,
    arc=0mm
]
\small
\textbf{Time Complexity:}
$O(n^3)$
\end{tcolorbox}


\begin{tcolorbox}[
    colback=violet!20,
    colframe=black,
    arc=0mm
]
\small
\textbf{Space Complexity:}
$O(n^2)$
\end{tcolorbox}


\begin{tcolorbox}[
    colback=pink!20,
    colframe=black,
    arc=0mm
]
\small
\textbf{Recurrence Formula:}
$m[i,j]=\min_{i \leq k < j}
(m[i,k]+m[k+1,j]+p[i-1]p[k]p[j])$
\end{tcolorbox}

\end{tcolorbox}
































\begin{tcolorbox}[
    colback=gray!30,
    colframe=black,
    boxrule=1pt,
    arc=0mm
]

\begin{tcolorbox}[
    colback=gray!10,
    colframe=black,
    arc=0mm
]
\centering
{\Huge\bfseries Longest Common Subsequence}
\end{tcolorbox}


\begin{tcolorbox}[
    colback=gray!10,
    colframe=black,
    arc=0mm
]
\small
Find the longest subsequence appearing in the same order in two sequences.
The algorithm uses dynamic programming by storing solutions of smaller
prefixes to build the optimal subsequence.
\end{tcolorbox}


\begin{tcolorbox}[
    colback=black!5,
    colframe=black,
    arc=0mm,
    title=Pseudo-code
]
\tiny
\begin{lstlisting}
LCS-LENGTH(X, Y, m, n)

    let b[1..m,1..n] and c[0..m,0..n]


    for i = 1 to m

        c[i,0] = 0


    for j = 0 to n

        c[0,j] = 0


    for i = 1 to m

        for j = 1 to n

            if x_i == y_j

                c[i,j] = c[i-1,j-1] + 1

                b[i,j] = "<-"


            else if c[i-1,j] >= c[i,j-1]

                c[i,j] = c[i-1,j]

                b[i,j] = "up"


            else

                c[i,j] = c[i,j-1]

                b[i,j] = "<-"


    return c and b
\end{lstlisting}
\end{tcolorbox}


\begin{tcolorbox}[
    colback=cyan!20,
    colframe=black,
    arc=0mm
]
\small
\textbf{Time Complexity:}
$O(mn)$
\end{tcolorbox}


\begin{tcolorbox}[
    colback=violet!20,
    colframe=black,
    arc=0mm
]
\small
\textbf{Space Complexity:}
$O(mn)$
\end{tcolorbox}


\begin{tcolorbox}[
    colback=pink!20,
    colframe=black,
    arc=0mm
]
\small
\textbf{Recurrence Formula:}

If $x_i=y_j$:
$c[i,j]=c[i-1,j-1]+1$

Else:
$c[i,j]=\max(c[i-1,j],c[i,j-1])$
\end{tcolorbox}

\end{tcolorbox}












\begin{tcolorbox}[
    colback=gray!30,
    colframe=black,
    boxrule=1pt,
    arc=0mm
]

\begin{tcolorbox}[
    colback=gray!10,
    colframe=black,
    arc=0mm
]
\centering
{\Huge\bfseries Print LCS}
\end{tcolorbox}


\begin{tcolorbox}[
    colback=gray!10,
    colframe=black,
    arc=0mm
]
\small
Reconstruct the longest common subsequence using the direction table
computed by the LCS dynamic programming algorithm.
The algorithm follows the stored arrows to rebuild the optimal sequence.
\end{tcolorbox}


\begin{tcolorbox}[
    colback=black!5,
    colframe=black,
    arc=0mm,
    title=Pseudo-code
]
\tiny
\begin{lstlisting}
PRINT-LCS(b, X, i, j)

    if i == 0 or j == 0

        return


    if b[i,j] == "diagonal"

        PRINT-LCS(b, X, i-1, j-1)

        print x_i


    elseif b[i,j] == "up"

        PRINT-LCS(b, X, i-1, j)


    else

        PRINT-LCS(b, X, i, j-1)
\end{lstlisting}
\end{tcolorbox}


\begin{tcolorbox}[
    colback=cyan!20,
    colframe=black,
    arc=0mm
]
\small
\textbf{Time Complexity:}
$O(m+n)$
\end{tcolorbox}


\begin{tcolorbox}[
    colback=violet!20,
    colframe=black,
    arc=0mm
]
\small
\textbf{Space Complexity:}
$O(m+n)$ recursion stack.
\end{tcolorbox}


\begin{tcolorbox}[
    colback=pink!20,
    colframe=black,
    arc=0mm
]
\small
\textbf{Recurrence Formula:}

$T(i,j)=T(i-1,j-1)+\Theta(1)$  
or $T(i,j)=T(i-1,j)+\Theta(1)$ depending on direction.
\end{tcolorbox}

\end{tcolorbox}















\begin{tcolorbox}[
    colback=gray!30,
    colframe=black,
    boxrule=1pt,
    arc=0mm
]

\begin{tcolorbox}[
    colback=gray!10,
    colframe=black,
    arc=0mm
]
\centering
{\Huge\bfseries Bottom-Up Optimal BST}
\end{tcolorbox}


\begin{tcolorbox}[
    colback=gray!10,
    colframe=black,
    arc=0mm
]
\small
Construct a binary search tree with minimum expected search cost.
The algorithm uses dynamic programming to find the optimal root
and recursively builds the best possible subtrees.
\end{tcolorbox}


\begin{tcolorbox}[
    colback=black!5,
    colframe=black,
    arc=0mm,
    title=Pseudo-code
]
\tiny
\begin{lstlisting}
OPTIMAL-BST(p, q, n)

    let e, w, root be new tables


    for i = 1 to n+1

        e[i,i-1] = 0

        w[i,i-1] = 0


    for l = 1 to n

        for i = 1 to n-l+1

            j = i+l-1

            e[i,j] = 8

            w[i,j] = w[i,j-1] + p[j]


            for r = i to j

                t = e[i,r-1]
                    + e[r+1,j]
                    + w[i,j]


                if t < e[i,j]

                    e[i,j] = t

                    root[i,j] = r


    return e and root
\end{lstlisting}
\end{tcolorbox}


\begin{tcolorbox}[
    colback=cyan!20,
    colframe=black,
    arc=0mm
]
\small
\textbf{Time Complexity:}
$O(n^3)$
\end{tcolorbox}


\begin{tcolorbox}[
    colback=violet!20,
    colframe=black,
    arc=0mm
]
\small
\textbf{Space Complexity:}
$O(n^2)$
\end{tcolorbox}


\begin{tcolorbox}[
    colback=pink!20,
    colframe=black,
    arc=0mm
]
\small
\textbf{Recurrence Formula:}

$e[i,j]=\min_{i \leq r \leq j}
(e[i,r-1]+e[r+1,j]+w[i,j])$
\end{tcolorbox}

\end{tcolorbox}














\begin{tcolorbox}[
    colback=gray!30,
    colframe=black,
    boxrule=1pt,
    arc=0mm
]

\begin{tcolorbox}[
    colback=gray!10,
    colframe=black,
    arc=0mm
]
\centering
{\Huge\bfseries Breadth-First Search}
\end{tcolorbox}


\begin{tcolorbox}[
    colback=gray!10,
    colframe=black,
    arc=0mm
]
\small
Explore a graph level by level starting from a source vertex $s$.
BFS computes the shortest distance in number of edges from $s$ to every
reachable vertex by using a queue.
\end{tcolorbox}


\begin{tcolorbox}[
    colback=black!5,
    colframe=black,
    arc=0mm,
    title=Pseudo-code
]
\tiny
\begin{lstlisting}
BFS(G, s)

    for each u in V - {s}

        u.d = 8


    s.d = 0

    Q = empty queue

    ENQUEUE(Q, s)


    while Q != empty

        u = DEQUEUE(Q)


        for each v in Adj[u]

            if v.d == 8

                v.d = u.d + 1

                ENQUEUE(Q, v)
\end{lstlisting}
\end{tcolorbox}


\begin{tcolorbox}[
    colback=cyan!20,
    colframe=black,
    arc=0mm
]
\small
\textbf{Time Complexity:}
$O(V+E)$
\end{tcolorbox}


\begin{tcolorbox}[
    colback=violet!20,
    colframe=black,
    arc=0mm
]
\small
\textbf{Space Complexity:}
$O(V)$
\end{tcolorbox}


\begin{tcolorbox}[
    colback=pink!20,
    colframe=black,
    arc=0mm
]
\small
\textbf{Recurrence Formula:}

$d[v]=d[u]+1$ when $v$ is discovered from $u$.
\end{tcolorbox}

\end{tcolorbox}














\begin{tcolorbox}[
    colback=gray!30,
    colframe=black,
    boxrule=1pt,
    arc=0mm
]

\begin{tcolorbox}[
    colback=gray!10,
    colframe=black,
    arc=0mm
]
\centering
{\Huge\bfseries Depth-First Search}
\end{tcolorbox}


\begin{tcolorbox}[
    colback=gray!10,
    colframe=black,
    arc=0mm
]
\small
Explore a graph by going as deep as possible along each branch before
backtracking. DFS assigns discovery and finishing times to vertices and
uses colors to track the exploration state.
\end{tcolorbox}


\begin{tcolorbox}[
    colback=black!5,
    colframe=black,
    arc=0mm,
    title=Pseudo-code
]
\tiny
\begin{lstlisting}
DFS(G)

    for each u in G.V

        u.color = WHITE


    time = 0


    for each u in G.V

        if u.color == WHITE

            DFS-VISIT(G,u)
\end{lstlisting}


\begin{lstlisting}
DFS-VISIT(G, u)

    time = time + 1

    u.d = time

    u.color = GRAY


    for each v in Adj[u]

        if v.color == WHITE

            DFS-VISIT(G,v)


    u.color = BLACK

    time = time + 1

    u.f = time
\end{lstlisting}
\end{tcolorbox}


\begin{tcolorbox}[
    colback=cyan!20,
    colframe=black,
    arc=0mm
]
\small
\textbf{Time Complexity:}
$O(V+E)$
\end{tcolorbox}


\begin{tcolorbox}[
    colback=violet!20,
    colframe=black,
    arc=0mm
]
\small
\textbf{Space Complexity:}
$O(V)$
\end{tcolorbox}


\begin{tcolorbox}[
    colback=pink!20,
    colframe=black,
    arc=0mm
]
\small
\textbf{Recurrence Formula:}

$T(V,E)=T(V-1,E)+O(E)$
$\Rightarrow O(V+E)$
\end{tcolorbox}

\end{tcolorbox}























\begin{tcolorbox}[
    colback=gray!30,
    colframe=black,
    boxrule=1pt,
    arc=0mm
]

\begin{tcolorbox}[
    colback=gray!10,
    colframe=black,
    arc=0mm
]
\centering
{\Huge\bfseries Topological Sort}
\end{tcolorbox}


\begin{tcolorbox}[
    colback=gray!10,
    colframe=black,
    arc=0mm
]
\small
Produce a linear ordering of the vertices of a Directed Acyclic Graph (DAG)
such that for every edge $(u,v)$, $u$ appears before $v$.
The algorithm uses DFS finishing times to determine the ordering.
\end{tcolorbox}


\begin{tcolorbox}[
    colback=black!5,
    colframe=black,
    arc=0mm,
    title=Pseudo-code
]
\tiny
\begin{lstlisting}
TOPOLOGICAL-SORT(G)

    DFS(G)

    // compute finishing times v.f


    output vertices in decreasing order

    of their finishing times
\end{lstlisting}
\end{tcolorbox}


\begin{tcolorbox}[
    colback=cyan!20,
    colframe=black,
    arc=0mm
]
\small
\textbf{Time Complexity:}
$O(V+E)$
\end{tcolorbox}


\begin{tcolorbox}[
    colback=violet!20,
    colframe=black,
    arc=0mm
]
\small
\textbf{Space Complexity:}
$O(V)$
\end{tcolorbox}


\begin{tcolorbox}[
    colback=pink!20,
    colframe=black,
    arc=0mm
]
\small
\textbf{Recurrence Formula:}

DFS traversal:
$T(V,E)=O(V+E)$
\end{tcolorbox}

\end{tcolorbox}













\begin{tcolorbox}[
    colback=gray!30,
    colframe=black,
    boxrule=1pt,
    arc=0mm
]

\begin{tcolorbox}[
    colback=gray!10,
    colframe=black,
    arc=0mm
]
\centering
{\Huge\bfseries Strongly Connected Components}
\end{tcolorbox}


\begin{tcolorbox}[
    colback=gray!10,
    colframe=black,
    arc=0mm
]
\small
Decompose a directed graph into maximal sets of vertices where every
vertex is reachable from every other vertex.
The algorithm uses two DFS passes and the transpose graph.
\end{tcolorbox}


\begin{tcolorbox}[
    colback=black!5,
    colframe=black,
    arc=0mm,
    title=Pseudo-code
]
\tiny
\begin{lstlisting}
SCC(G)

    DFS(G)

    // compute finishing times u.f


    Compute G^T

    // reverse all edges


    DFS(G^T)

    considering vertices

    in decreasing order of u.f


    output each DFS tree

    as one SCC
\end{lstlisting}
\end{tcolorbox}


\begin{tcolorbox}[
    colback=cyan!20,
    colframe=black,
    arc=0mm
]
\small
\textbf{Time Complexity:}
$O(V+E)$
\end{tcolorbox}


\begin{tcolorbox}[
    colback=violet!20,
    colframe=black,
    arc=0mm
]
\small
\textbf{Space Complexity:}
$O(V+E)$
\end{tcolorbox}


\begin{tcolorbox}[
    colback=pink!20,
    colframe=black,
    arc=0mm
]
\small
\textbf{Recurrence Formula:}

Two DFS traversals:
$T(V,E)=2O(V+E)=O(V+E)$
\end{tcolorbox}

\end{tcolorbox}














\begin{tcolorbox}[
    colback=gray!30,
    colframe=black,
    boxrule=1pt,
    arc=0mm
]

\begin{tcolorbox}[
    colback=gray!10,
    colframe=black,
    arc=0mm
]
\centering
{\Huge\bfseries Maximum Flow - Ford Fulkerson}
\end{tcolorbox}


\begin{tcolorbox}[
    colback=gray!10,
    colframe=black,
    arc=0mm
]
\small
Find the maximum possible flow from a source $s$ to a sink $t$ in a flow
network while respecting capacities.
The algorithm repeatedly finds augmenting paths in the residual graph and
increases the flow until no improvement is possible.
\end{tcolorbox}


\begin{tcolorbox}[
    colback=black!5,
    colframe=black,
    arc=0mm,
    title=Pseudo-code
]
\tiny
\begin{lstlisting}
FORD-FULKERSON(G, s, t)

    initialize flow f = 0


    while there exists an

    augmenting path p in Gf


        augment flow f along p


        update residual graph


    return f
\end{lstlisting}
\end{tcolorbox}


\begin{tcolorbox}[
    colback=cyan!20,
    colframe=black,
    arc=0mm
]
\small
\textbf{Time Complexity:}

$O(E \cdot |f^*|)$

for integer capacities.
\end{tcolorbox}


\begin{tcolorbox}[
    colback=violet!20,
    colframe=black,
    arc=0mm
]
\small
\textbf{Space Complexity:}

$O(V+E)$
\end{tcolorbox}


\begin{tcolorbox}[
    colback=pink!20,
    colframe=black,
    arc=0mm
]
\small
\textbf{Recurrence Formula:}

$f = f + f_p$

where $f_p$ is the bottleneck capacity of the augmenting path.
\end{tcolorbox}

\end{tcolorbox}
























\begin{tcolorbox}[
    colback=gray!30,
    colframe=black,
    boxrule=1pt,
    arc=0mm
]

\begin{tcolorbox}[
    colback=gray!10,
    colframe=black,
    arc=0mm
]
\centering
{\Huge\bfseries Connected Components}
\end{tcolorbox}


\begin{tcolorbox}[
    colback=gray!10,
    colframe=black,
    arc=0mm
]
\small
Find all connected components of an undirected graph using the
Union-Find data structure. Each vertex starts in its own set, and edges
merge sets that belong to the same component.
\end{tcolorbox}


\begin{tcolorbox}[
    colback=black!5,
    colframe=black,
    arc=0mm,
    title=Pseudo-code
]
\tiny
\begin{lstlisting}
CONNECTED-COMPONENTS(G)

    for each vertex v in G.V

        MAKE-SET(v)


    for each edge (u,v) in G.E

        if FIND-SET(u) != FIND-SET(v)

            UNION(u,v)
\end{lstlisting}
\end{tcolorbox}


\begin{tcolorbox}[
    colback=cyan!20,
    colframe=black,
    arc=0mm
]
\small
\textbf{Time Complexity:}

$O(E \cdot \alpha(V))$
\end{tcolorbox}


\begin{tcolorbox}[
    colback=violet!20,
    colframe=black,
    arc=0mm
]
\small
\textbf{Space Complexity:}

$O(V)$
\end{tcolorbox}


\begin{tcolorbox}[
    colback=pink!20,
    colframe=black,
    arc=0mm
]
\small
\textbf{Recurrence Formula:}

Union-Find operations:
$O(\alpha(V))$ amortized per operation.
\end{tcolorbox}

\end{tcolorbox}














\begin{tcolorbox}[
    colback=gray!30,
    colframe=black,
    boxrule=1pt,
    arc=0mm
]

\begin{tcolorbox}[
    colback=gray!10,
    colframe=black,
    arc=0mm
]
\centering
{\Huge\bfseries MAKE-SET}
\end{tcolorbox}


\begin{tcolorbox}[
    colback=gray!10,
    colframe=black,
    arc=0mm
]
\small
Create a new disjoint set containing only one element.
Each element starts as its own root in the Union-Find structure.
\end{tcolorbox}


\begin{tcolorbox}[
    colback=black!5,
    colframe=black,
    arc=0mm,
    title=Pseudo-code
]
\tiny
\begin{lstlisting}
MAKE-SET(x)

    x.p = x

    x.rank = 0
\end{lstlisting}
\end{tcolorbox}


\begin{tcolorbox}[
    colback=cyan!20,
    colframe=black,
    arc=0mm
]
\small
\textbf{Time Complexity:}

$O(1)$
\end{tcolorbox}


\begin{tcolorbox}[
    colback=violet!20,
    colframe=black,
    arc=0mm
]
\small
\textbf{Space Complexity:}

$O(1)$ per element
\end{tcolorbox}


\begin{tcolorbox}[
    colback=pink!20,
    colframe=black,
    arc=0mm
]
\small
\textbf{Recurrence Formula:}

Single-node tree initialization.
\end{tcolorbox}

\end{tcolorbox}












\begin{tcolorbox}[
    colback=gray!30,
    colframe=black,
    boxrule=1pt,
    arc=0mm
]

\begin{tcolorbox}[
    colback=gray!10,
    colframe=black,
    arc=0mm
]
\centering
{\Huge\bfseries FIND-SET}
\end{tcolorbox}


\begin{tcolorbox}[
    colback=gray!10,
    colframe=black,
    arc=0mm
]
\small
Find the representative (root) of the set containing an element.
The operation uses path compression to flatten the tree and make future
queries faster.
\end{tcolorbox}


\begin{tcolorbox}[
    colback=black!5,
    colframe=black,
    arc=0mm,
    title=Pseudo-code
]
\tiny
\begin{lstlisting}
FIND-SET(x)

    if x != x.p

        x.p = FIND-SET(x.p)


    return x.p
\end{lstlisting}
\end{tcolorbox}


\begin{tcolorbox}[
    colback=cyan!20,
    colframe=black,
    arc=0mm
]
\small
\textbf{Time Complexity:}

$O(\alpha(n))$ amortized
\end{tcolorbox}


\begin{tcolorbox}[
    colback=violet!20,
    colframe=black,
    arc=0mm
]
\small
\textbf{Space Complexity:}

$O(1)$ auxiliary
\end{tcolorbox}


\begin{tcolorbox}[
    colback=pink!20,
    colframe=black,
    arc=0mm
]
\small
\textbf{Recurrence Formula:}

Path compression:
$T(n) \approx O(\alpha(n))$
\end{tcolorbox}

\end{tcolorbox}















\begin{tcolorbox}[
    colback=gray!30,
    colframe=black,
    boxrule=1pt,
    arc=0mm
]

\begin{tcolorbox}[
    colback=gray!10,
    colframe=black,
    arc=0mm
]
\centering
{\Huge\bfseries UNION}
\end{tcolorbox}


\begin{tcolorbox}[
    colback=gray!10,
    colframe=black,
    arc=0mm
]
\small
Merge two disjoint sets into one set.
The operation finds the representatives of both sets and links the two
trees together using the Union-Find structure.
\end{tcolorbox}


\begin{tcolorbox}[
    colback=black!5,
    colframe=black,
    arc=0mm,
    title=Pseudo-code
]
\tiny
\begin{lstlisting}
UNION(x, y)

    LINK(FIND-SET(x),
         FIND-SET(y))
\end{lstlisting}
\end{tcolorbox}


\begin{tcolorbox}[
    colback=cyan!20,
    colframe=black,
    arc=0mm
]
\small
\textbf{Time Complexity:}

$O(\alpha(n))$ amortized
\end{tcolorbox}


\begin{tcolorbox}[
    colback=violet!20,
    colframe=black,
    arc=0mm
]
\small
\textbf{Space Complexity:}

$O(1)$ auxiliary
\end{tcolorbox}


\begin{tcolorbox}[
    colback=pink!20,
    colframe=black,
    arc=0mm
]
\small
\textbf{Recurrence Formula:}

Union by rank:
$T(n) \approx O(\alpha(n))$
\end{tcolorbox}

\end{tcolorbox}





















\begin{tcolorbox}[
    colback=gray!30,
    colframe=black,
    boxrule=1pt,
    arc=0mm
]

\begin{tcolorbox}[
    colback=gray!10,
    colframe=black,
    arc=0mm
]
\centering
{\Huge\bfseries LINK}
\end{tcolorbox}


\begin{tcolorbox}[
    colback=gray!10,
    colframe=black,
    arc=0mm
]
\small
Attach one root under another in the Union-Find structure.
The operation uses union by rank to keep the trees shallow and maintain
an efficient balanced forest.
\end{tcolorbox}


\begin{tcolorbox}[
    colback=black!5,
    colframe=black,
    arc=0mm,
    title=Pseudo-code
]
\tiny
\begin{lstlisting}
LINK(x, y)

    if x.rank > y.rank

        y.p = x


    else

        x.p = y


        if x.rank == y.rank

            y.rank = y.rank + 1
\end{lstlisting}
\end{tcolorbox}


\begin{tcolorbox}[
    colback=cyan!20,
    colframe=black,
    arc=0mm
]
\small
\textbf{Time Complexity:}

$O(1)$
\end{tcolorbox}


\begin{tcolorbox}[
    colback=violet!20,
    colframe=black,
    arc=0mm
]
\small
\textbf{Space Complexity:}

$O(1)$ auxiliary
\end{tcolorbox}


\begin{tcolorbox}[
    colback=pink!20,
    colframe=black,
    arc=0mm
]
\small
\textbf{Recurrence Formula:}

Union by rank:
keeps tree height bounded by $O(\log n)$.
\end{tcolorbox}

\end{tcolorbox}

















\begin{tcolorbox}[
    colback=gray!30,
    colframe=black,
    boxrule=1pt,
    arc=0mm
]

\begin{tcolorbox}[
    colback=gray!10,
    colframe=black,
    arc=0mm
]
\centering
{\Huge\bfseries PRIM'S ALGORITHM}
\end{tcolorbox}


\begin{tcolorbox}[
    colback=gray!10,
    colframe=black,
    arc=0mm
]
\small
Prim's algorithm builds a Minimum Spanning Tree (MST) of a weighted graph.
It repeatedly chooses the minimum weight edge connecting the current tree
to a new vertex using a priority queue.
\end{tcolorbox}


\begin{tcolorbox}[
    colback=black!5,
    colframe=black,
    arc=0mm,
    title=Pseudo-code
]
\tiny
\begin{lstlisting}
PRIM(G, w, r)

    Q = empty priority queue

    for each u in G.V
        u.key = 8
        u.pi = NIL
        INSERT(Q, u)

    DECREASE-KEY(Q, r, 0)

    while Q != empty

        u = EXTRACT-MIN(Q)

        for each v in G.Adj[u]

            if v in Q and w(u,v) < v.key

                v.pi = u
                DECREASE-KEY(Q, v, w(u,v))
\end{lstlisting}
\end{tcolorbox}


\begin{tcolorbox}[
    colback=cyan!20,
    colframe=black,
    arc=0mm
]
\small
\textbf{Time Complexity:}

$O(E \log V)$ with binary heap

$O(E + V \log V)$ with Fibonacci heap
\end{tcolorbox}


\begin{tcolorbox}[
    colback=violet!20,
    colframe=black,
    arc=0mm
]
\small
\textbf{Space Complexity:}

$O(V)$
\end{tcolorbox}


\begin{tcolorbox}[
    colback=pink!20,
    colframe=black,
    arc=0mm
]
\small
\textbf{Recurrence Formula:}

Greedy expansion:

$T(V,E)=O(E\log V)$
\end{tcolorbox}

\end{tcolorbox}
































\begin{tcolorbox}[
    colback=gray!30,
    colframe=black,
    boxrule=1pt,
    arc=0mm
]

\begin{tcolorbox}[
    colback=gray!10,
    colframe=black,
    arc=0mm
]
\centering
{\Huge\bfseries KRUSKAL'S ALGORITHM}
\end{tcolorbox}


\begin{tcolorbox}[
    colback=gray!10,
    colframe=black,
    arc=0mm
]
\small
Kruskal's algorithm constructs a Minimum Spanning Tree (MST) by selecting
the smallest weight edges while avoiding cycles using the Union-Find data structure.
\end{tcolorbox}


\begin{tcolorbox}[
    colback=black!5,
    colframe=black,
    arc=0mm,
    title=Pseudo-code
]
\tiny
\begin{lstlisting}
KRUSKAL(G, w)

    A = empty set

    for each vertex v in G.V

        MAKE-SET(v)


    sort edges of G.E by increasing weight


    for each edge (u,v) in sorted order

        if FIND-SET(u) != FIND-SET(v)

            A = A U {(u,v)}
            UNION(u,v)


    return A
\end{lstlisting}
\end{tcolorbox}


\begin{tcolorbox}[
    colback=cyan!20,
    colframe=black,
    arc=0mm
]
\small
\textbf{Time Complexity:}

$O(E\log E)$

$=O(E\log V)$
\end{tcolorbox}


\begin{tcolorbox}[
    colback=violet!20,
    colframe=black,
    arc=0mm
]
\small
\textbf{Space Complexity:}

$O(V)$
\end{tcolorbox}


\begin{tcolorbox}[
    colback=pink!20,
    colframe=black,
    arc=0mm
]
\small
\textbf{Recurrence Formula:}

Sorting dominates:

$T(E)=O(E\log E)$
\end{tcolorbox}

\end{tcolorbox}























\begin{tcolorbox}[
    colback=gray!30,
    colframe=black,
    boxrule=1pt,
    arc=0mm
]

\begin{tcolorbox}[
    colback=gray!10,
    colframe=black,
    arc=0mm
]
\centering
{\Huge\bfseries RELAX}
\end{tcolorbox}


\begin{tcolorbox}[
    colback=gray!10,
    colframe=black,
    arc=0mm
]
\small
Relaxation improves the shortest path estimate of a vertex by checking
whether passing through another vertex gives a shorter route.
\end{tcolorbox}


\begin{tcolorbox}[
    colback=black!5,
    colframe=black,
    arc=0mm,
    title=Pseudo-code
]
\tiny
\begin{lstlisting}
RELAX(u, v, w)

    if v.d > u.d + w(u,v)

        v.d = u.d + w(u,v)
        v.pi = u
\end{lstlisting}
\end{tcolorbox}


\begin{tcolorbox}[
    colback=cyan!20,
    colframe=black,
    arc=0mm
]
\small
\textbf{Time Complexity:}

$O(1)$
\end{tcolorbox}


\begin{tcolorbox}[
    colback=violet!20,
    colframe=black,
    arc=0mm
]
\small
\textbf{Space Complexity:}

$O(1)$
\end{tcolorbox}


\begin{tcolorbox}[
    colback=pink!20,
    colframe=black,
    arc=0mm
]
\small
\textbf{Recurrence Formula:}

Shortest path update:

$d[v]=\min(d[v],d[u]+w(u,v))$
\end{tcolorbox}

\end{tcolorbox}















\begin{tcolorbox}[
    colback=gray!30,
    colframe=black,
    boxrule=1pt,
    arc=0mm
]

\begin{tcolorbox}[
    colback=gray!10,
    colframe=black,
    arc=0mm
]
\centering
{\Huge\bfseries BELLMAN-FORD}
\end{tcolorbox}


\begin{tcolorbox}[
    colback=gray!10,
    colframe=black,
    arc=0mm
]
\small
Bellman-Ford computes shortest paths from a source vertex and supports
negative edge weights. It repeatedly relaxes all edges to improve estimates.
\end{tcolorbox}


\begin{tcolorbox}[
    colback=black!5,
    colframe=black,
    arc=0mm,
    title=Pseudo-code
]
\tiny
\begin{lstlisting}
BELLMAN-FORD(G,w,s)

    INIT-SINGLE-SOURCE(G,s)

    for i = 1 to |G.V|-1

        for each edge (u,v) in G.E

            RELAX(u,v,w)


    return TRUE
\end{lstlisting}
\end{tcolorbox}


\begin{tcolorbox}[
    colback=cyan!20,
    colframe=black,
    arc=0mm
]
\small
\textbf{Time Complexity:}

$O(VE)$
\end{tcolorbox}


\begin{tcolorbox}[
    colback=violet!20,
    colframe=black,
    arc=0mm
]
\small
\textbf{Space Complexity:}

$O(V)$
\end{tcolorbox}


\begin{tcolorbox}[
    colback=pink!20,
    colframe=black,
    arc=0mm
]
\small
\textbf{Recurrence Formula:}

$V-1$ relaxations:

$T(V,E)=O(VE)$
\end{tcolorbox}

\end{tcolorbox}






















\begin{tcolorbox}[
    colback=gray!30,
    colframe=black,
    boxrule=1pt,
    arc=0mm
]

\begin{tcolorbox}[
    colback=gray!10,
    colframe=black,
    arc=0mm
]
\centering
{\Huge\bfseries DIJKSTRA}
\end{tcolorbox}


\begin{tcolorbox}[
    colback=gray!10,
    colframe=black,
    arc=0mm
]
\small
Dijkstra's algorithm computes shortest paths from a source vertex in a graph
with non-negative edge weights. It repeatedly selects the closest unprocessed
vertex and relaxes its outgoing edges.
\end{tcolorbox}


\begin{tcolorbox}[
    colback=black!5,
    colframe=black,
    arc=0mm,
    title=Pseudo-code
]
\tiny
\begin{lstlisting}
DIJKSTRA(G,w,s)

    INIT-SINGLE-SOURCE(G,s)

    S = empty set
    Q = G.V

    while Q != empty

        u = EXTRACT-MIN(Q)

        S = S U {u}


        for each v in G.Adj[u]

            RELAX(u,v,w)
\end{lstlisting}
\end{tcolorbox}


\begin{tcolorbox}[
    colback=cyan!20,
    colframe=black,
    arc=0mm
]
\small
\textbf{Time Complexity:}

$O(E\log V)$ with binary heap

$O(V\log V + E)$ with efficient decrease-key
\end{tcolorbox}


\begin{tcolorbox}[
    colback=violet!20,
    colframe=black,
    arc=0mm
]
\small
\textbf{Space Complexity:}

$O(V)$
\end{tcolorbox}


\begin{tcolorbox}[
    colback=pink!20,
    colframe=black,
    arc=0mm
]
\small
\textbf{Recurrence Formula:}

Greedy shortest path:

$T(V,E)=O((V+E)\log V)$
\end{tcolorbox}

\end{tcolorbox}


















































\end{document}